\documentclass[12pt,a4paper]{article}
\usepackage[utf8]{inputenc}
\usepackage[T1]{fontenc}
\usepackage[french]{babel}
\usepackage{amsmath, amssymb}
\usepackage{geometry}
\usepackage{graphicx}
\usepackage{xcolor}
\usepackage{hyperref}
\usepackage{enumitem}
\usepackage{fancyhdr}
\usepackage{titlesec}
\usepackage{mdframed}
\usepackage{tcolorbox}
\tcbuselibrary{skins, breakable}

\geometry{margin=2.5cm}
\setlength{\headheight}{14.5pt}

% En-tête et pied de page
\pagestyle{fancy}
\fancyhf{}
\lhead{BTS -- Mathématiques}
\rhead{TP Info -- Séries de Fourier}
\cfoot{\thepage}

% ---- Couleurs ----
\definecolor{exobg}{RGB}{232,240,254}
\definecolor{exoframe}{RGB}{70,130,220}
\definecolor{rappelbg}{RGB}{255,243,220}
\definecolor{rappelframe}{RGB}{230,150,30}
\definecolor{repbg}{RGB}{245,255,245}
\definecolor{repframe}{RGB}{80,170,80}
\definecolor{fichierbg}{RGB}{240,240,255}
\definecolor{fichierframe}{RGB}{100,100,200}

% ---- Boîte : rappel théorique ----
\newmdenv[
    backgroundcolor=rappelbg,
    linecolor=rappelframe,
    linewidth=1.5pt,
    roundcorner=5pt,
    innertopmargin=8pt,
    innerbottommargin=8pt,
]{rappelbox}

% ---- Boîte : exercice / question ----
\newmdenv[
    backgroundcolor=exobg,
    linecolor=exoframe,
    linewidth=1.5pt,
    roundcorner=5pt,
    innertopmargin=8pt,
    innerbottommargin=8pt,
]{exercicebox}

% ---- Boîte : fichier Python à lancer ----
\newmdenv[
    backgroundcolor=fichierbg,
    linecolor=fichierframe,
    linewidth=1.5pt,
    roundcorner=5pt,
    innertopmargin=6pt,
    innerbottommargin=6pt,
]{fichierbox}

% ---- Boîte réponse (lignes vides pour l'élève) ----
\newcommand{\cadreReponse}[1][4]{%
    \begin{mdframed}[
        backgroundcolor=repbg,
        linecolor=repframe,
        linewidth=1pt,
        roundcorner=4pt,
        innertopmargin=6pt,
        innerbottommargin=6pt,
    ]
    \vspace{#1\baselineskip}
    \end{mdframed}%
}

% ---- Formatage des titres (12pt, numérotation avec points) ----
\titleformat{\section}[block]{\normalsize\bfseries}{\thesection.}{0.5em}{}
\titleformat{\subsection}[block]{\normalsize\bfseries}{\thesubsection.}{0.5em}{}

% Pour \foreach
\usepackage{pgffor}

% ---------------------------------------------------------------
\begin{document}

\thispagestyle{fancy}

\begin{center}
    {\LARGE \textbf{TP Info -- Séries de Fourier 2}}\\[4pt]
    {\large Calcul des coefficients et tracé des approximations}\\[8pt]
    \rule{0.6\linewidth}{0.4pt}\\[6pt]
    {\small Durée : 2 heures \quad -- \quad Documents autorisés : cours, aide Python}
\end{center}

\medskip
\noindent\textbf{Nom :} \dotfill \quad \textbf{Prénom :} \dotfill \quad \textbf{Classe :} \makebox[3cm]{\dotfill}

\bigskip

\hrule

% ---------------------------------------------------------------
\section*{Objectifs}

Ce TP a pour but de mettre en œuvre le calcul numérique des coefficients de Fourier d'une fonction périodique, puis de visualiser les approximations partielles de la série de Fourier associée.

À l'issue de ce TP, vous serez capables de :
\begin{itemize}
    \item Calculer numériquement les coefficients $a_0$, $a_n$ et $b_n$ d'une série de Fourier ;
    \item Reconstruire l'approximation partielle d'ordre $N$ d'une fonction ;
    \item Visualiser le phénomène de Gibbs ;
    \item Mesurer l'erreur d'approximation.
\end{itemize}

\begin{fichierbox}
\textbf{Fichiers Python fournis :} ouvrez le dossier du TP, vous y trouverez les scripts suivants à lancer dans l'ordre des exercices.\\[4pt]
\begin{tabular}{ll}
    \texttt{ex1\_creneau.py}  & Exercice 1 -- Fonction créneau \\
    \texttt{ex2\_triangle.py} & Exercice 2 -- Fonction triangle \\
    \texttt{ex3\_erreur.py}   & Exercice 3 -- Erreur d'approximation \\
    \texttt{ex4\_libre.py}    & Exercice 4 -- Fonction au choix (à modifier) \\
\end{tabular}\\[4pt]
Pour lancer un script : ouvrir un terminal dans le dossier du TP, puis taper \texttt{python ex1\_creneau.py}
\end{fichierbox}

% ---------------------------------------------------------------
\section{Rappels théoriques}

\begin{rappelbox}
Soit $f$ une fonction $T$-périodique, intégrable sur une période. Sa \textbf{série de Fourier} est :
\[
    f(x) \;\approx\; \frac{a_0}{2} + \sum_{n=1}^{+\infty} \left[ a_n \cos\!\left(\frac{2\pi n x}{T}\right) + b_n \sin\!\left(\frac{2\pi n x}{T}\right) \right]
\]
avec les \textbf{coefficients de Fourier} :
\[
    a_0 = \frac{2}{T} \int_{0}^{T} f(x)\, dx
    \qquad
    a_n = \frac{2}{T} \int_{0}^{T} f(x)\cos\!\left(\frac{2\pi n x}{T}\right) dx \]
    
    \[b_n = \frac{2}{T} \int_{0}^{T} f(x)\sin\!\left(\frac{2\pi n x}{T}\right) dx \]

\end{rappelbox}

\medskip
L'\textbf{approximation partielle} d'ordre $N$ est :
\[
    S_N(x) = \frac{a_0}{2} + \sum_{n=1}^{N} \left[ a_n \cos\!\left(\frac{2\pi n x}{T}\right) + b_n \sin\!\left(\frac{2\pi n x}{T}\right) \right]
\]

% ---------------------------------------------------------------
\section{Exercice 1 -- Fonction créneau \hfill\normalsize\texttt{ex1\_creneau.py}}

\begin{fichierbox}
Lancez le script \texttt{ex1\_creneau.py}. Il trace la fonction créneau, calcule les coefficients $a_0$, $a_n$, $b_n$ pour $n=1,\ldots,10$, puis affiche les approximations partielles $S_N$ pour $N \in \{1, 3, 5, 15, 50\}$.
\end{fichierbox}

On s'intéresse à la \textbf{fonction créneau} $T$-périodique avec $T = 2\pi$ :
\[
    f(x) = \begin{cases} 1 & \text{si } x \in [0, \pi[ \\ -1 & \text{si } x \in [\pi, 2\pi[ \end{cases}
\]

\subsection{Observation de la fonction}

\begin{exercicebox}
\textbf{Question 1.1 :} Observez le premier graphique affiché par le script.
\end{exercicebox}

\medskip
\noindent\textbf{a)} La fonction créneau est-elle paire, impaire, ou ni l'une ni l'autre ?

\cadreReponse[1]

\noindent\textbf{b)} Qu'en déduisez-vous sur les coefficients $a_n$ ($n \geq 1$) ?

\cadreReponse[1]
\fbox{fermer la fenêtre du graphique. } 

\subsection{Coefficients de Fourier} 

\begin{exercicebox}
\textbf{Question 1.2 :} Analysez le tableau de coefficients affiché dans le terminal.
\end{exercicebox}

\medskip
\noindent\textbf{a)} Les valeurs numériques de $a_n$ confirment-elles votre réponse à la question 1.1 ?

\cadreReponse[2]

\noindent\textbf{b)} Recopiez les cinq premières valeurs non nulles de $b_n$ dans le tableau ci-dessous.

\medskip
\begin{center}
\renewcommand{\arraystretch}{1.6}
\begin{tabular}{|c|c|c|c|c|c|}
    \hline
    $n$ & 1 & 2 & 3 & 4 & 5 \\
    \hline
    $b_n$ & \hspace{2cm} & \hspace{2cm} & \hspace{2cm} & \hspace{2cm} & \hspace{2cm} \\
    \hline
\end{tabular}
\end{center}

\medskip
\noindent\textbf{c)} Pour quelles valeurs de $n$ les coefficients $b_n$ sont-ils nuls ?

\cadreReponse[1]

\newpage

\subsection{Approximations partielles et phénomène de Gibbs}

\begin{exercicebox}
\textbf{Question 1.3 :} Observez les graphiques des approximations $S_N$.
\end{exercicebox}

\vfill 

\noindent\textbf{a)} Que constatez-vous aux points de discontinuité ($-2\pi$, $-\pi$, 0, $\pi$ et $2\pi$ )quand $N$ augmente ? 

\cadreReponse[3]

\vfill  


\noindent\textbf{b)} Modifiez la liste \texttt{ordres} dans le script pour ajouter $N = 100$. Le phénomène disparaît-il complètement ?

\cadreReponse[2]

\vfill  


\begin{rappelbox}
\textbf{Le phénomène de Gibbs.}
Lorsqu'on approche par une série de Fourier une fonction présentant une \textbf{discontinuité de saut} (comme la fonction créneau), il apparaît des oscillations parasites de part et d'autre du saut : c'est le \textbf{phénomène de Gibbs}.

Ce phénomène a deux caractéristiques remarquables :
\begin{itemize}
    \item Quand $N$ augmente, les oscillations se \textbf{resserrent} autour du point de discontinuité, mais leur amplitude \textbf{ne diminue pas} : elle reste voisine de $9\,\%$ de la hauteur du saut.
    \item Le phénomène \textbf{ne disparaît jamais}, quelle que soit la valeur de $N$ : c'est une limite fondamentale des séries de Fourier pour les fonctions discontinues.
\end{itemize}

En revanche, pour une fonction \textbf{continue} (comme la fonction triangle), ce phénomène n'existe pas et la convergence est bien meilleure.
\end{rappelbox}

% ---------------------------------------------------------------

\vfill 

\section{Exercice 2 -- Fonction triangle \hfill\normalsize\texttt{ex2\_triangle.py}}

\begin{fichierbox}
Lancez le script \texttt{ex2\_triangle.py}. Il calcule les coefficients et trace les approximations de la fonction triangle.
\end{fichierbox}

\vfill 

On étudie la \textbf{fonction triangle} $T$-périodique avec $T = 2\pi$ :
\[
    f(x) = \left| x - \pi \right| - \frac{\pi}{2}, \quad x \in [0, 2\pi[
\]

\vfill 

\newpage 


\begin{exercicebox}
\textbf{Question 2.1 :} Parité et coefficients.
\end{exercicebox}

\medskip
\noindent\textbf{a)} La fonction triangle est-elle paire, impaire, ou ni l'une ni l'autre ? Qu'en déduisez-vous sur les $b_n$ ?

\cadreReponse[2]

\noindent\textbf{b)} Vérifiez sur le tableau affiché que $a_0 \approx 0$ et que les $b_n$ sont bien nuls.

\cadreReponse[2]

\begin{exercicebox}
\textbf{Question 2.2 :} Comparaison de la convergence.
\end{exercicebox}

\medskip
\noindent\textbf{a)} Comparez les graphiques $S_N$ pour $N = 1, 5, 20$ avec ceux de l'exercice 1 au même ordre. Laquelle des deux fonctions est mieux approchée à ordre $N$ égal ?

\cadreReponse[3]

\noindent\textbf{b)} La fonction triangle est-elle continue ? Dérivable ? Quel lien peut-on faire avec la vitesse de convergence ?

\cadreReponse[3]

% ---------------------------------------------------------------
\section{Exercice 3 -- Mesure de l'erreur d'approximation \hfill\normalsize\texttt{ex3\_erreur.py}}

\begin{fichierbox}
Lancez le script \texttt{ex3\_erreur.py}. Il calcule et trace l'erreur quadratique moyenne entre $f$ et $S_N$ en fonction de $N$, pour les deux fonctions étudiées. \textbf{Attention, le calcul peut prendre quelques secondes.}
\end{fichierbox}

\medskip
L'\textbf{erreur quadratique} d'ordre $N$ est définie par :
\[
    \varepsilon_N = \sqrt{\frac{1}{P} \sum_{k=1}^{P} \left( f(x_k) - S_N(x_k) \right)^2}
\]
où $x_1, \ldots, x_P$ est une discrétisation de la période.

\begin{exercicebox}
\textbf{Question 3.1 :} Analysez les courbes de convergence (graphique en échelle logarithmique).
\end{exercicebox}

\medskip
\noindent\textbf{a)} Quelle fonction converge le plus rapidement ? Cela confirme-t-il vos observations de l'exercice 2 ?

\cadreReponse[2]

\noindent\textbf{b)} Les courbes semblent-elles linéaires sur le graphique en échelle log ? Qu'indique une droite sur ce type de graphique ?

\cadreReponse[2]

\noindent\textbf{c)} D'après le tableau affiché dans le terminal, à partir de quel ordre $N$ l'erreur sur le créneau passe-t-elle en dessous de $0.1$ ?

\cadreReponse[2]

% ---------------------------------------------------------------
\section{Exercice 4 -- Fonction au choix \hfill\normalsize\texttt{ex4\_libre.py}}

\begin{fichierbox}
Ouvrez le fichier \texttt{ex4\_libre.py} dans un éditeur. Modifiez la fonction \texttt{f\_libre} en remplaçant le corps par la formule de votre choix ci-dessous, puis sauvegardez et lancez le script.
\end{fichierbox}

\medskip
Choisissez \textbf{une} des fonctions suivantes (cochez votre choix) :

\begin{center}
\renewcommand{\arraystretch}{1.8}
\begin{tabular}{|c|l|l|}
    \hline
    \textbf{Choix} & \textbf{Nom} & \textbf{Formule} \\
    \hline
    $\square$ & Dent de scie    & $f(x) = \dfrac{x \bmod T}{T} - \dfrac{1}{2}$ \\
    \hline
    $\square$ & Gaussienne périodique & $f(x) = e^{-\cos(x)}$ \\
    \hline
    $\square$ & Parabolique     & $f(x) = (x - \pi)^2 - \dfrac{\pi^2}{3}$ \quad pour $x \in [0, 2\pi[$ \\
    \hline
\end{tabular}
\end{center}

\newpage

\begin{exercicebox}
\textbf{Question 4.1 :} Pour la fonction choisie, répondez aux questions suivantes après avoir lancé le script.
\end{exercicebox}

\medskip
\noindent\textbf{a)} Notez les valeurs de $a_0$, $a_1$, $b_1$ affichées.

\cadreReponse[2]

\noindent\textbf{b)} La fonction est-elle paire, impaire ? Cela se confirme-t-il sur les coefficients ?

\cadreReponse[2]

\noindent\textbf{c)} La convergence est-elle plus rapide que pour le créneau ? Justifiez par la régularité de la fonction.

\cadreReponse[3]

% ---------------------------------------------------------------
\section*{Bilan}

\begin{exercicebox}
\textbf{Question de synthèse :} Complétez le tableau récapitulatif suivant.

\medskip
\begin{center}
\renewcommand{\arraystretch}{2}
\begin{tabular}{|l|c|c|c|c|}
    \hline
    \textbf{Fonction} & \textbf{Continue ?} & \textbf{Paire / Impaire} & \textbf{$b_n = 0$ ?} & \textbf{Convergence} \\
    \hline
    Créneau  & & & & \\
    \hline
    Triangle & & & & \\
    \hline
    Votre choix & & & & \\
    \hline
\end{tabular}
\end{center}

\medskip
Quelle relation générale peut-on établir entre la \textbf{régularité} d'une fonction et la \textbf{vitesse de décroissance} de ses coefficients de Fourier ?
\end{exercicebox}

\cadreReponse[5]

\vspace{1cm}
\hrule
\vspace{0.3cm}
\begin{center}
\textit{Bon courage !}
\end{center}

\end{document}
